\documentclass[10pt,a4paper]{amsart}
\usepackage[margin=19mm]{geometry}
\usepackage{amsmath,amssymb,mathtools}
\usepackage[scr=rsfs,cal=euler]{mathalfa}
\usepackage{xfrac}
\usepackage[unicode,hidelinks,bookmarks=false]{hyperref}
\newcommand{\ie}{i.e.\ }
\newcommand{\cf}{cf.\ }
\newcommand{\resp}{respectively\ }
\newcommand{\Subsection}{\subsection}
\providecommand{\nobreakdash}{\nobreak\hbox{-}}
\setlength{\emergencystretch}{2em}
\multlinegap0pt
\usepackage{bm}
\usepackage{bbm}
\usepackage{dsfont}
\usepackage{xspace}
\usepackage{cancel}
\usepackage{mathtools}
\usepackage{amsthm}
\makeatletter
\@ifundefined{proposition}{\newtheorem{proposition}{Proposition}}{}
\@ifundefined{lemma}{\newtheorem{lemma}[proposition]{Lemma}}{}
\makeatother
\DeclareMathOperator{\GL}{GL}
\DeclareMathOperator{\Mat}{Mat}
\DeclareMathOperator{\Sp}{Sp}
\newcommand{\Z}{\mathbb{Z}}
\title[Polarizations, torsors and theta groups]{Polarizations, torsors and theta groups}
\author{Jef Laga}
\date{Source: arXiv:2510.24678v2 (CC BY 4.0)}
\begin{document}
\maketitle

\noindent\emph{Source and licence.} Polarizations, torsors and theta groups. Version arXiv:2510.24678v2 (CC BY 4.0), distributed under the Creative Commons Attribution 4.0 International licence, as recorded in the arXiv licence metadata for this exact version and shown on the abstract page https://arxiv.org/abs/2510.24678v2. Full rights notice: licenses/arxiv-2510.24678-CC-BY-4.0.txt.\par\smallskip

\noindent\textit{Editorial scope.} Counted unit: the Lemma whose source label is \texttt{lemma: block diagonal elements are in the commutator subgroup} in the article (source0.tex lines 2584--2586, source label \texttt{lemma: block diagonal elements are in the commutator subgroup}), delivered together with the article's own complete original proof of it (source0.tex lines 2587--2661, one \texttt{proof} environment of 75 lines). No mathematics is written, repaired or re-derived: statement and proof are the article's own text line for line. Required context retained uncounted: lemma: abelianization level subgroups symplectic group (lines 2564--2571). Every definition, display and label of those lines is retained unchanged. Omitted from this package and not redistributed: 1--2563, 2572--2583, 2662--3156. Nothing inside a retained excerpt is omitted or altered: every retained body is delivered exactly as the article's lines are written, so no omission receipt of any kind accompanies this package. Citations: the retained excerpts carry no citation command at all, so no bibliography entry of the article is transcribed and no reference is redistributed. Reference adaptation: none was needed and none was made; every reference inside the delivered unit resolves inside it, so no unresolved reference or citation is produced. Printed slips kept literal: no printed word, symbol, hypothesis or number of the counted statement or of its proof was altered. Layout and class: the article's own class is \texttt{article} with packages including amsmath, amssymb, amsthm, fontenc, inputenc, enumitem, graphicx, geometry, hyperref, tikz, tikz-cd, xcolor, stmaryrd, mathrsfs; the wrapper is the lane's pinned \texttt{amsart} editorial wrapper at 10pt on A4, which restates the theorem-like environments the retained text needs and adds the article's own macro definitions that the retained text needs (\texttt{\textbackslash GL}, \texttt{\textbackslash Mat}, \texttt{\textbackslash Sp}, \texttt{\textbackslash Z}), with extra wrapper packages (bm, bbm, dsfont, xspace, cancel, mathtools). The delivered numbering is the unit's own. Assets: no retained span contains a figure environment, a graphics-inclusion command, a PDF-page inclusion, a table or a TikZ picture; no image file is copied into this package, the package declares no asset dependency and no image file is redistributed with it. Licence: version arXiv:2510.24678v2 of this article is distributed under the Creative Commons Attribution 4.0 International licence, as recorded in the arXiv licence metadata for this exact version and shown on the abstract page https://arxiv.org/abs/2510.24678v2. A full rights notice is delivered at licenses/arxiv-2510.24678-CC-BY-4.0.txt. Characters: every retained excerpt is pure ASCII, so the delivered engine, XeLaTeX with its default Latin Modern text font, reports no missing character.
\begin{lemma}\label{lemma: abelianization level subgroups symplectic group}
    If $n\geq 3$, $\Sp_{2n}(\Z_2)^{\mathrm{ab}}$ is trivial.
    If $k\geq 2$, the commutator subgroup of $\Sp_{2k}(\Z_2,2)$ equals 
    \[
    \Sp_{2k}(\Z_2,4,8):=
    \{ A = (a_{ij}) \in \Sp_{2k}(\Z_2)\colon A\equiv I \bmod 4 \text{  and }a_{i,i+k}\equiv a_{i+k,i}\equiv 0 \bmod{8} \text{ for all }i=1,\dots, k\}.
    \]
\end{lemma}

\begin{lemma}\label{lemma: block diagonal elements are in the commutator subgroup}
    If $n\geq 3$ and $k\geq 2$, then $L\subset \Gamma^{\mathrm{der}}$.
\end{lemma}

\begin{proof}
    By considering block diagonal elements of $\Sp_{2n}(\Z_2)$, we see that every element of $L$ of the form 
    \begin{align}\label{equation: first commutator lemma, first type element}
    \alpha\left(\begin{pmatrix}
        X_{11} & 0 \\ 0 & I 
    \end{pmatrix}\right), \quad X_{11}\in \GL_n(\Z_2)
    \end{align}
    lies in the image of $i\colon \Sp_{2n}(\Z_2)\rightarrow \Gamma$. 
    By Lemma \ref{lemma: abelianization level subgroups symplectic group} and the assumption $n\geq 3$, every element in this image lies in $\Gamma^{\mathrm{der}}$, so every element of the form \eqref{equation: first commutator lemma, first type element} lies in $\Gamma^{\mathrm{der}}$.
    Similarly, every element of the form 
    \begin{align}\label{equation: first commutator lemma, second type element}
        \alpha\left( 
        \begin{pmatrix}
            I & 0 \\ 0 & X_{22} 
        \end{pmatrix}
        \right),\quad 
        X_{22}\in \ker(\GL_k(\Z_2)\rightarrow \GL_k(\Z/4\Z))
    \end{align}
    lies in $\Gamma^{\mathrm{der}}$.
    Next, we consider strictly upper triangular elements.
    If $X\in \GL_n(\Z_2)$ and $Y\in 2\mathrm{Mat}_{n,k}(\Z_2)$, we calculate the commutator
    \[
    \left[
    \begin{pmatrix}X & 0 \\ 0 & I \end{pmatrix},
    \begin{pmatrix}I & Y \\ 0 & I \end{pmatrix}
    \right]
    =
    \begin{pmatrix}I & XY-Y \\ 0 & I \end{pmatrix}
    \]
    in $L'$.
    We claim that the $\Z$-span of the set $\{XY-Y\colon X\in \GL_n(\Z_2), Y\in 2\mathrm{Mat}_{n,k}(\Z_2) \}$ equals $2\mathrm{Mat}_{n,k}(\Z_2)$. 
    To prove this, we may assume $Y$ consists of a single column, in which case it can be explicitly checked by taking $X$ to be upper and lower triangular matrices (and using that $n\geq 2$).
    We conclude that every element of the form $\alpha\left( \begin{psmallmatrix}
        I & * \\ 0 & I 
    \end{psmallmatrix}\right) \in L$ lies in $L^{\mathrm{der}}$.
    By an identical argument, every element of the form $\alpha\left( \begin{psmallmatrix}
        I & 0 \\ * & I 
    \end{psmallmatrix}\right) \in L$ lies in $L^{\mathrm{der}}$.
    Now let $\alpha\left(\begin{psmallmatrix}X_{11} & X_{12} \\ X_{21} & X_{22}  \end{psmallmatrix}\right)\in L$ be a general element.
    Then the condition $X_{12} \equiv 0 \bmod 2$ implies that $X_{11}$ is invertible modulo $2$, so $X_{11}\in \GL_n(\Z_2)$. 
    Therefore, left multiplying by $\alpha\left(
    \begin{psmallmatrix}
        I & 0 \\ -X_{21} & I
    \end{psmallmatrix}
    \begin{psmallmatrix}
        X_{11}^{-1} & 0 \\ 0 & I
    \end{psmallmatrix}
    \right)\in \Gamma^{\mathrm{der}}$ and right multiplying by $\alpha\left(
    \begin{psmallmatrix}
        I & -X_{11}^{-1} X_{12} \\ 0 & I
    \end{psmallmatrix}
    \right)\in \Gamma^{\mathrm{der}}$ shows that
    \begin{align}\label{equation: first commutator lemma, reduction to lower block matrix}
    \alpha\left(\begin{pmatrix}X_{11} & X_{12} \\ X_{21} & X_{22}  \end{pmatrix} \right)\equiv
    \alpha\left(\begin{pmatrix}I & 0 \\ 0 & X_{22}-X_{21}X_{11}^{-1}X_{12}  \end{pmatrix}\right) 
    \mod \Gamma^{\mathrm{der}}.
    \end{align}
    To prove the lemma, it therefore suffices to show that every element of $L$ the form $\alpha\left( \begin{psmallmatrix}
        I & 0 \\ 0 & * 
    \end{psmallmatrix}\right)$ lies in $\Gamma^{\mathrm{der}}$.
    We calculate using \eqref{equation: first commutator lemma, reduction to lower block matrix} that for all $X\in \Mat_{n,k}(\Z_2)$ and $Y\in \Mat_{k,n}(\Z_2)$:
    \[
    \alpha\left(\left[ 
    \begin{pmatrix}I & 2X \\ 0 & I \end{pmatrix},
    \begin{pmatrix}I & 0 \\ Y & I \end{pmatrix}
    \right]\right)=
    \alpha\left(\begin{pmatrix}* & -4XYX \\ 2YXY & I-2YX \end{pmatrix}\right) \equiv 
    \alpha\left(\begin{pmatrix}I & 0 \\ 0 & I-2YX+ 8(\cdots) \end{pmatrix} \right)
    \mod \Gamma^{\mathrm{der}}
    ,
    \]
    where $(\cdots)$ denotes an element of $\Mat_{k,k}(\Z_2)$.
    By combining this identity with the elements of the form \eqref{equation: first commutator lemma, second type element}, it suffices to prove that the set $\{I-2YX\colon X\in \Mat_{n,k}(\Z_2), \,Y\in\Mat_{k,n}(\Z_2)\}$ generates a subgroup of $\GL_k(\Z_2)$ that surjects onto $\ker(\GL_k(\Z/4\Z)\rightarrow \GL_k(\Z/2\Z))$.  
    This is true, as can be seen by taking $X$ and $Y$ to be $\Z_2$-multiples of elementary matrices.
\end{proof}

\end{document}
